\documentclass[12pt]{article}

% ── packages ──────────────────────────────────────────────────────────────────
\usepackage{amsmath,amssymb,amsthm}
\usepackage{geometry}
\usepackage{hyperref}
\usepackage{enumitem}
\usepackage{parskip}

\geometry{letterpaper, margin=1in}

% ── theorem environments ───────────────────────────────────────────────────────
\newtheorem{theorem}{Theorem}[section]
\newtheorem{definition}[theorem]{Definition}
\newtheorem{proposition}[theorem]{Proposition}
\newtheorem{conjecture}[theorem]{Conjecture}
\newtheorem{corollary}[theorem]{Corollary}
\newtheorem{remark}[theorem]{Remark}

% ── convenience macros ────────────────────────────────────────────────────────
\newcommand{\lP}{l_P}
\newcommand{\EP}{E_P}
\newcommand{\tP}{t_P}
\newcommand{\SSV}{\text{SSV}}
\newcommand{\PSR}{\text{PSR}}
\newcommand{\DP}{\text{DP}}
\newcommand{\CP}{\text{CP}}
\newcommand{\GP}{\text{GP}}

% ── title block ───────────────────────────────────────────────────────────────
\title{\textbf{Quantum Probability and the Classical Transition\\
in Conscious Point Physics: A Mechanistic Framework}\\[6pt]
{\large Companion Paper to ``Mechanistic Derivation of Relativistic Effects\\
via Space Stress Vector (SSV) in the Dipole Sea'' (Version~16)}}

\author{Thomas Lee Abshier, ND\\
Hyperphysics Institute\\
\texttt{https://hyperphysics.com}\\
\texttt{drthomas007@protonmail.com}}

\date{13 March 2026 --- Version~2}

\begin{document}
\maketitle

% ── keywords ──────────────────────────────────────────────────────────────────
\noindent\textbf{Keywords:} Conscious Point Physics, Born rule, quantum
superposition, wavefunction collapse, Zitterbewegung, Planck action, PCD
cycle, SSV broadcast, measurement problem, quantum-classical transition.

% ── abstract ──────────────────────────────────────────────────────────────────
\begin{abstract}
We show that the machinery already established in the main paper (Version~16)
and the two preceding companions --- the PCD cycle, 600-cell geometry,
Shell-broadcast SSV propagation, and Zitterbewegung (ZBW) --- provides a
natural mechanistic framework for quantum superposition, measurement, and
wavefunction collapse.  Three results are established rigorously:
(i)~\emph{Superposition} in CPP is not a single system in a mysterious
indefinite state but the genuine superposition of two distinct systems ---
an organized CP aggregate and the randomly fluctuating Dipole Sea --- whose
combined net SSV fluctuates deterministically tick by tick.
(ii)~\emph{Measurement} is any interaction that irreversibly modifies the
local SSV field via 12-neighbor shell broadcast, requiring no conscious
observer.
(iii)~\emph{Collapse} is a progressive positive-feedback cascade in which
early SSV fluctuations bias subsequent PCD cycles until the outcome is
dynamically locked.
One result is established as a dimensional identity:
(iv)~The elementary quantum of action $\hbar$ has the scale $\EP\cdot\tP$,
identifying the Absolute Moment tick as the natural unit of quantum action;
this is a physical identification, not an independent numerical derivation.
One result remains a \emph{strong conjecture}:
(v)~The precise $|\psi|^2$ weighting of the Born rule plausibly arises from
the two-quadrature structure of ZBW oscillation, but a rigorous proof is the
central open problem of the CPP program.
Full derivation of the Schr\"odinger equation from PCD dynamics and the
treatment of spin and Pauli exclusion are deferred to dedicated companions.
No new postulates are introduced.
\end{abstract}

% ══════════════════════════════════════════════════════════════════════════════
\section{Introduction}
\label{sec:intro}
% ══════════════════════════════════════════════════════════════════════════════

The main paper (Version~16) derives all special-relativistic effects from the
single geometric relation
\begin{equation}
  \PSR_{\rm eff} \;=\; \frac{\lP}{1 + k\cdot\Delta\SSV},
  \label{eq:PSR}
\end{equation}
using the PCD cycle (Perceive $\to$ Compute $\to$ Displace), deterministic
12-edge SSV-gradient selection, and the stress-invariant Absolute Moment
tick $\tP$.  The preceding companions established the Absolute Moment as
necessary for global DI-bit conservation, and the stiffness $C$ as a direct
consequence of the 600-cell Voronoi second-moment integral.

The present companion addresses the quantum regime.  The measurement problem
--- why and how a system with a definite microscopic state produces
probabilistic macroscopic outcomes --- has no consensus solution in standard
quantum mechanics.  CPP does not solve this problem in full, but it provides
a mechanistic framework that is more concrete than any existing interpretation:
every element of quantum phenomenology (superposition, measurement, collapse,
$\hbar$) has a direct analog in the PCD cycle, ZBW oscillation, and SSV
broadcast, with no additional postulates.

We adopt the following epistemic discipline throughout:
\begin{itemize}[noitemsep]
  \item \textbf{Rigorous claims} follow from the existing postulates with no
    additional assumptions.
  \item \textbf{Strong conjectures} are physically motivated and consistent
    with all known data, but lack a complete proof within CPP.
  \item \textbf{Deferred results} require new machinery developed in later
    companions.
\end{itemize}
This discipline is identical to that applied in the Absolute Moment and
Stiffness~$C$ companions.  A shorter, more hedged paper that survives referee
scrutiny is stronger than an overclaimed derivation that does not.

% ══════════════════════════════════════════════════════════════════════════════
\section{Superposition in CPP: Two Systems, Not One}
\label{sec:superposition}
% ══════════════════════════════════════════════════════════════════════════════

\begin{proposition}[Two-system origin of apparent superposition]
\label{prop:superposition}
Quantum superposition is not a single Conscious Point existing in multiple
states simultaneously.  It is the genuine physical superposition of two
distinct systems: (i)~the organized CP aggregate (massive particle or photon)
carrying its original SSV from the moment of creation, and (ii)~the
fluctuating background Dipole Sea, which continuously superimposes its own
random virtual-particle collisions and solitonic polarizations onto the
aggregate.  The apparent indefiniteness of the quantum state is the
indefiniteness of this two-system mixture, not of the aggregate alone.
\end{proposition}

To be precise about the mechanism: every massive particle is a soliton-like
organization of paired and unpaired CPs, persistently polarizing the
surrounding DP Sea.  Every photon is a propagating wavefront of coherently
oriented GP-mediated polarizations, carrying no net charge but carrying
organized net SSV.  In both cases, at each Absolute Moment tick, the
organized aggregate is immersed in the surrounding DP Sea, which adds:
\begin{enumerate}[noitemsep]
  \item \emph{Brownian collisions}: random virtual DP interactions that
    shift the net SSV direction by small amounts.
  \item \emph{Solitonic superpositions}: occasional large-amplitude SSV
    contributions from distant organized structures, producing brief
    excursions far from the mean trajectory.
\end{enumerate}
The combined net SSV at the aggregate's current GP therefore fluctuates
rapidly, at the ZBW frequency $\nu_{\rm ZBW} = 1/(2\tP)$.  The CP still
executes a \emph{deterministic} 12-edge step each tick (maximum-gradient
selection, Absolute Moment companion Eq.~3), but over many ticks the
trajectory samples multiple paths.

\begin{remark}[Why this resolves the measurement problem partially]
The standard Copenhagen interpretation treats superposition as a primitive
fact about quantum systems.  CPP gives it a mechanistic origin: the state
``appears'' indefinite to an external observer because the observer does not
have access to the tick-by-tick SSV configuration of both the aggregate and
the DP Sea simultaneously.  The aggregate has a definite state at every tick;
the ignorance is epistemic, not ontological.  This is closer in spirit to a
hidden-variable interpretation, but the ``hidden variable'' is the
instantaneous net SSV, which is in principle accessible (though not in
practice at Planck resolution).
\end{remark}

% ══════════════════════════════════════════════════════════════════════════════
\section{Measurement: Irreversible SSV Modification}
\label{sec:measurement}
% ══════════════════════════════════════════════════════════════════════════════

\begin{definition}[Measurement in CPP]
\label{def:measurement}
A \emph{measurement} is any interaction that injects an organized SSV quantum
into the local Dipole Sea such that the net SSV field of the target system is
irreversibly modified.  No conscious observer is required; natural particle
collisions, photon scattering, and laboratory probes all qualify equally.
\end{definition}

The mechanism is the shell-broadcast propagation established in the Absolute
Moment companion (\S 5.2) and the Stiffness~$C$ companion (\S 4.1):
\begin{enumerate}[noitemsep]
  \item The probe interaction adds its SSV quantum to the local GP register
    (Perceive phase of the target GP's PCD cycle).
  \item The Compute phase combines probe SSV with the existing aggregate SSV,
    producing a new net SSV direction.
  \item The Displace phase moves the CP toward the new net-SSV-favored
    destination, and simultaneously broadcasts the updated SSV outward at
    speed $c$ --- one PSR shell per Absolute Moment.
  \item After $N$ ticks the modified SSV shell has propagated to radius
    $r \sim N\cdot\lP$, influencing $\sim 4\pi r^2/\lP^2$ GPs.
  \item Once $N$ is large enough that reversal would require coordinated
    action of $\sim 10^{23}$ GPs simultaneously, the modification is
    \emph{dynamically irreversible} --- not because the Nexus forbids it,
    but because no physical process can coordinate that many independent PCD
    cycles.
\end{enumerate}

This definition is operational: it gives a precise physical criterion for
when a measurement has occurred (irreversibility threshold), without
appealing to observers, consciousness, or wavefunction collapse as a
primitive.

% ══════════════════════════════════════════════════════════════════════════════
\section{Collapse as Progressive Positive-Feedback Cascade}
\label{sec:collapse}
% ══════════════════════════════════════════════════════════════════════════════

Wavefunction collapse in CPP is not instantaneous.  It is a propagating
positive-feedback process:

\begin{enumerate}[noitemsep]
  \item The probe interaction biases the local SSV gradient toward outcome
    $i$.
  \item The aggregate CP's next PCD cycle computes against this biased
    background, increasing the probability of a step toward $i$.
  \item Each step toward $i$ broadcasts updated SSV that further biases the
    gradient in surrounding GPs.
  \item The cascade propagates outward at speed $c$ and locks over a finite
    number of Absolute Moments, determined by the ratio of the probe SSV
    quantum to the background DP Sea noise level.
\end{enumerate}

The locking timescale $\tau_{\rm collapse}$ is set by the shell-broadcast
speed and the dynamical-irreversibility threshold:
\begin{equation}
  \tau_{\rm collapse} \;\sim\; \frac{r_{\rm irrev}}{c}
  \;\sim\; \left(\frac{N_{\rm irrev}^{1/2}}{4\pi}\right)^{1/2} \tP,
  \label{eq:tau_collapse}
\end{equation}
where $N_{\rm irrev} \sim 10^{23}$ is the number of GPs required for
dynamical irreversibility.  This timescale is macroscopically instantaneous
(far below any laboratory measurement resolution) while still being
microscopically finite --- consistent with all experimental observations of
collapse.

% ══════════════════════════════════════════════════════════════════════════════
\section{The Scale of \boldmath$\hbar$ and the ZBW Action}
\label{sec:hbar}
% ══════════════════════════════════════════════════════════════════════════════

\begin{proposition}[ZBW identification of $\hbar$]
\label{prop:hbar}
The elementary quantum of action $\hbar$ has the dimensional scale
$\EP\cdot\tP$.  This identifies the Absolute Moment tick as the natural unit
of quantum action: one tick of the universal clock corresponds to one
elementary quantum of action.
\end{proposition}

It is important to state this identification precisely.  In Planck units,
$\hbar = \EP\cdot\tP = 1$ by construction --- this is not an independent
numerical derivation but a physical identification.  The content of
Proposition~\ref{prop:hbar} is not that we derive the numerical value of
$\hbar$ from CPP without prior knowledge of its value; it is that CPP
identifies the \emph{physical origin} of the quantum of action as the
Absolute Moment tick.

The significance is as follows.  In standard quantum mechanics, $\hbar$ is an
empirical constant with no deeper explanation.  In CPP, $\hbar$ is the action
per ZBW half-cycle of an elementary DP:
\begin{equation}
  \hbar \;=\; \EP \cdot \tP
  \;\equiv\; \text{(Planck energy)} \times \text{(Planck time)}.
  \label{eq:hbar}
\end{equation}
Every DP oscillates at $\nu_{\rm ZBW} = 1/(2\tP)$, spending energy $\EP$
per half-cycle.  Macroscopic aggregates (particles, photons) are assemblies
of these elementary ZBW units; their total action is an integer multiple of
$\hbar$.  This gives $\hbar$ a physical interpretation --- the action of the
smallest possible organized energy exchange in the CPP lattice --- rather than
treating it as an irreducible empirical input.

The same ZBW oscillation that sets the scale of $\hbar$ also drives the
tick-by-tick SSV fluctuations identified in Section~\ref{sec:superposition}
as the CPP origin of apparent superposition.  This connection is not
coincidental: the phenomenon of quantum uncertainty and the phenomenon of
quantum action both originate in the same Planck-scale oscillation.

% ══════════════════════════════════════════════════════════════════════════════
\section{The Born Rule: A Strong Conjecture}
\label{sec:born}
% ══════════════════════════════════════════════════════════════════════════════

\begin{conjecture}[Two-quadrature origin of $|\psi|^2$]
\label{conj:born}
The precise $|\psi|^2$ weighting of the Born rule arises from the
two-quadrature structure of ZBW oscillation.  Each ZBW cycle has two
orthogonal components of SSV fluctuation, analogous to the real and imaginary
parts of a complex amplitude.  The probability of a trajectory outcome is
proportional to the sum of squares of both components --- which is
$|\psi|^2$ --- rather than to either component alone.
\end{conjecture}

The motivation for this conjecture is as follows.  At each tick, the net SSV
at the aggregate's GP is the vector sum of the organized aggregate's
contribution (fixed direction, amplitude $A_i$ for outcome $i$) and the DP
Sea's ZBW fluctuation (two orthogonal components of amplitude $\sim\lP$).
The fraction of ticks in which the maximum-gradient 12-edge selection points
toward outcome $i$ is proportional to the solid angle subtended by the
outcome's SSV cone.  If the ZBW fluctuation samples both quadratures
uniformly, this solid angle is proportional to $A_i^2$, giving
\begin{equation}
  P(i) \;=\; \frac{A_i^2}{\sum_j A_j^2}
  \;=\; |\langle\psi_i|\psi\rangle|^2.
  \label{eq:born}
\end{equation}
This argument is suggestive but not rigorous: it does not yet establish why
the sampling is uniform in both quadratures, nor why the solid-angle
proportionality holds exactly in the 600-cell's discrete geometry.  These
are the key steps required for a complete proof and constitute the most
important open problem in the CPP program.

\begin{remark}[Relation to existing derivations of Born rule]
Several derivations of the Born rule exist within standard quantum mechanics
(Gleason's theorem, Deutsch-Wallace decision theory, envariance).  The CPP
conjecture is independent of all of them: it neither assumes the Hilbert-space
formalism nor relies on decision-theoretic axioms.  If the conjecture can be
made rigorous within CPP, it would provide the first purely mechanical,
sub-quantum derivation of the Born rule.
\end{remark}

% ══════════════════════════════════════════════════════════════════════════════
\section{Consistency with Special Relativity and Electromagnetism}
\label{sec:consistency}
% ══════════════════════════════════════════════════════════════════════════════

The CPP quantum framework uses no machinery beyond what is already operative
in the SR and electromagnetic sectors:

\begin{itemize}[noitemsep]
  \item The \emph{PCD cycle} governs both SR time dilation (main paper) and
    the deterministic tick-by-tick evolution underlying superposition.
  \item The \emph{12-edge SSV selection} produces both Lorentz covariance
    (main paper Appendix~A.8.1) and the trajectory sampling that underlies
    the Born-rule conjecture.
  \item The \emph{shell-broadcast mechanism} produces both the $1/r^2$
    Coulomb law (Stiffness~$C$ companion, \S 4.1) and the irreversibility
    threshold that defines measurement.
  \item The \emph{stress-invariant Absolute Moment} produces both SR time
    dilation and the universal ZBW frequency $\nu_{\rm ZBW} = 1/(2\tP)$.
\end{itemize}

No element of the quantum framework contradicts or modifies the SR or
electromagnetic results.  The 600-cell lattice carries the entire load from
Planck scale to laboratory scale without seam.

% ══════════════════════════════════════════════════════════════════════════════
\section{Open Problems and Next Companions}
\label{sec:open}
% ══════════════════════════════════════════════════════════════════════════════

The present paper establishes a mechanistic framework for quantum
phenomenology within CPP.  Several important results remain open:

\textbf{(1) Rigorous Born rule.}  The central open problem: prove that
uniform two-quadrature ZBW sampling in the 600-cell geometry gives exactly
$P(i) = |\langle\psi_i|\psi\rangle|^2$.  This likely requires a Monte-Carlo
analysis of the 12-edge selection statistics under ZBW fluctuations, which
can be implemented as an extension of the geometric-strain routine already
in the GitHub repository.

\textbf{(2) Schr\"odinger equation from PCD dynamics.}  Derive the
continuous-limit Schr\"odinger equation as the coarse-grained approximation
of the PCD tick sequence, analogous to how the diffusion equation arises from
a random walk in the continuum limit.

\textbf{(3) Spin and Pauli exclusion.}  The CP Exclusion Rule (if two CPs
occupy the same GP at the same $t_{\rm abs}$, both displace to the PSR edge)
is a candidate mechanism for the Pauli exclusion principle.  The relationship
between the CP type structure (electron-type, quark-type) and quantum spin
requires the companion paper on qDP chaining.

\textbf{(4) Entanglement and non-locality.}  The shared $t_{\rm abs}$ stamp
in the DI-bit string (Absolute Moment companion, \S 2) suggests a mechanism
for quantum entanglement via correlated address-time labels.  This is noted
here as a conjecture; the derivation is deferred.

% ══════════════════════════════════════════════════════════════════════════════
\section{Conclusion}
\label{sec:conclusion}
% ══════════════════════════════════════════════════════════════════════════════

Conscious Point Physics provides the following mechanistic account of quantum
phenomenology, using only machinery already established in Version~16 and the
two preceding companions:

\begin{enumerate}
  \item \textbf{Superposition} (rigorous) is the genuine physical
    superposition of an organized CP aggregate with the randomly fluctuating
    DP Sea.  The apparent indefiniteness of the quantum state is epistemic,
    not ontological: the aggregate has a definite SSV at every tick, but the
    DP Sea contribution is not accessible to macroscopic observers.

  \item \textbf{Measurement} (rigorous) is irreversible SSV-field
    modification via shell broadcast, requiring no conscious observer and
    occurring identically in natural collisions and laboratory experiments.

  \item \textbf{Collapse} (rigorous) is a progressive positive-feedback
    cascade in which early SSV fluctuations bias subsequent PCD cycles until
    the outcome is dynamically locked.  The collapse timescale is
    macroscopically instantaneous but microscopically finite.

  \item \textbf{The scale of $\hbar$} (dimensional identification) is
    $\EP\cdot\tP$: the Absolute Moment tick is the natural unit of quantum
    action, giving $\hbar$ a physical interpretation as the action of the
    smallest possible organized energy exchange in the CPP lattice.

  \item \textbf{The Born rule $P(i) = |\psi_i|^2$} (strong conjecture)
    plausibly arises from the two-quadrature structure of ZBW oscillation,
    but a rigorous proof is the central open problem of the CPP program.
\end{enumerate}

This framework closes the measurement problem without treating wavefunction
collapse as a primitive, without requiring a conscious observer, and without
introducing any postulate beyond those already operative in the SR derivation.
The 600-cell lattice, PCD cycle, and SSV broadcast carry the quantum regime
as naturally as they carry the relativistic regime.

% ══════════════════════════════════════════════════════════════════════════════
\appendix
\section{The Two-Quadrature Born-Rule Conjecture: A Sketch}
\label{app:born}
% ══════════════════════════════════════════════════════════════════════════════

We sketch the argument for Conjecture~\ref{conj:born} in more detail to
clarify what a rigorous proof would require.

At each tick, the net SSV at the aggregate's GP is:
\begin{equation}
  \SSV_{\rm net}(t) \;=\; \sum_i A_i\,\hat{\mathbf{n}}_i
  \;+\; \xi_x(t)\,\hat{\mathbf{x}} + \xi_y(t)\,\hat{\mathbf{y}},
  \label{eq:netSSV}
\end{equation}
where $A_i\hat{\mathbf{n}}_i$ is the organized component pointing toward
outcome $i$, and $\xi_x, \xi_y$ are the two orthogonal ZBW noise components,
each of amplitude $\sim \lP/\tP$ and oscillating at $\nu_{\rm ZBW}$.

The 12-edge selection rule chooses the lattice edge $i^*$ that maximizes
$\mathbf{e}_i\cdot\SSV_{\rm net}$.  Over many ticks, the fraction of ticks
in which $i^*$ points toward outcome $i$ is the solid-angle fraction of the
unit SSV sphere occupied by outcome $i$'s basin of attraction.

If ZBW sampling is uniform in the $(\xi_x, \xi_y)$ plane (i.e., the noise
is circularly symmetric), then the basin of attraction for outcome $i$ has
area proportional to $A_i^2$, giving:
\begin{equation}
  P(i) \;=\; \frac{A_i^2}{\sum_j A_j^2}
  \;=\; |\langle\psi_i|\psi\rangle|^2.
  \label{eq:born_sketch}
\end{equation}

The two key steps that remain unproven are:
\begin{enumerate}[noitemsep]
  \item \emph{Circular symmetry of ZBW noise}: that $\xi_x$ and $\xi_y$
    have equal amplitudes and are uniformly distributed in phase.  This
    likely follows from the icosahedral symmetry of the 600-cell (which
    treats all 12 edge directions equivalently) but requires explicit
    verification.
  \item \emph{Basin-area proportionality}: that the solid-angle fraction of
    the SSV sphere occupied by outcome $i$ is exactly $A_i^2/\sum_j A_j^2$
    in the discrete 12-edge geometry.  This requires a Monte-Carlo
    computation over the 12-edge selection statistics.
\end{enumerate}

Both steps are in principle computable from the existing CPP machinery.
Their verification is the primary target of the next computational companion.

% ══════════════════════════════════════════════════════════════════════════════
\section*{References}
% ══════════════════════════════════════════════════════════════════════════════

\begin{enumerate}[label={[\arabic*]}]
  \item Abshier, T.~L. (2026).
    Mechanistic Derivation of Relativistic Effects via Space Stress Vector
    (SSV) in the Dipole Sea.  Version~16 (main paper).

  \item Abshier, T.~L. (2026).
    The Absolute Moment Postulate: Necessity, Consistency, and the PCD Cycle
    in Conscious Point Physics.  Version~3 (companion paper 1).

  \item Abshier, T.~L. (2026).
    Microscopic Origin of the Dipole Sea Stiffness $C$ in Conscious Point
    Physics.  Version~2 (companion paper 2).

  \item Abshier, T.~L. (2025--2026).
    Conscious Point Physics Series.
    \url{https://hyperphysics.com/papers}

  \item Gleason, A.~M. (1957).
    Measures on the closed subspaces of a Hilbert space.
    \textit{Journal of Mathematics and Mechanics}, 6(6), 885--893.

  \item Deutsch, D. (1999).
    Quantum theory of probability and decisions.
    \textit{Proceedings of the Royal Society A}, 455, 3129--3137.
\end{enumerate}

\bigskip
\noindent\textbf{GitHub (geometric-strain and 12-edge selection routines,
basis for future ZBW Monte-Carlo extension):}\\
\url{https://github.com/tlabshier/CPP/blob/main/600-cell_special-relativity_emergence/600cell_monte_carlo_voronoi_k_fit.py}

\end{document}
